\section{Theory}
\label{sec-theory}

A model's answer is pulled toward the interest that is present in its
context. On an open-ended dilemma the interest is the asker's own account
of the conflict, on a false premise it is the premise the user owns, and
on an assigned side it is the side. The pull cannot be removed from inside
the same context, because the copy that would remove it is the copy that
carries it. Three copies of one model vote as one voter
and the gain from a correlated jury is bounded by how far its members'
errors are independent
\citep{ladha1992condorcet,boland1989majority,dietrich2013independent,kim2025correlated}.
What can be done is to place an opposed interest into the process so that
agreement across interests, or dissent against interest, becomes
informative. This is the argumentative theory of reasoning
\citep{mercier2011argumentative,mercier2017enigma} joined to the social
account of objectivity \citep{longino1990science,longino2002fate,popper2011open}.
Individuals argue for a side, and truth is a property of the structure of
criticism among them rather than of any one mind. The scaffold and the
collective are the two ways of placing the opposed interest.

Internal opposition is the scaffold. One trace is made to represent the
other parties. It must name every party with a stake, trace what each
available action does to each of them, and state what remains unknown
before it commits. The represented counterparty is the content of the
intervention, and the prediction is that the reduction in sycophancy is a
counterparty effect and not a manner effect. Three consequences are
testable. The reduction should be carried by the sections that represent
the other parties, the stakeholders and the consequences, rather than by
the section that licenses hedging, which is the competing account in our
own record. It should be concentrated on accounts that contain a party
other than the asker and absent where none exists, so that a false
theorem, which has no second party, is the case where the scaffold has
nothing to represent and where we already observe that it affirms the
falsehood more often rather than less.
And it should be a property of the prompt's structure, reproducing across
vendors and on open-weight models.

External opposition is the collective. When the interests are separate
seats the credible-signal rule applies. A statement against the speaker's
assigned interest is informative and a statement with it is not
\citep{milgrom1986relying,lipman1995robust,glazer2001debates}. A verdict
with one losing party gives that party's seat a reason to object, and its
objection carries nothing. A second objection has no side to explain it
and must come from the account. The dissent edge exists only where the
interests are opposed, it survives removing the exchange between seats and
removing narration from the seats, so it is the assignment and not the
conversation that carries it
\citep{dewatripont1999advocates,krishna2001model}.
The collective is a sensor and not a decider. No arrangement beats its best
seat by vote, the moderator's competence sets the verdict
\citep{shin1998adversarial},
and the gain is realised by routing the flagged cases to a stronger judge
with no side. The rule presupposes an interest, so it should fail on a
model whose seats do not hold their sides, and it does.
That is the precise sense in which the collective depends on the same
control the scaffold exploits. Embodiment has to hold for opposition to be
readable.

The game-theoretic reading is used only where its primitives exist, an
assigned interest, a receiver whose verdict is separate from the message,
and a decision read independently of the speaker. Those exist in the
collective and not in a single trace, so for the scaffold alone the paper
claims a change in the output distribution whose content is the
represented counterparty and makes no signalling claim. Naming what another
party knows and wants is also an expression-side theory-of-mind act, which
is why the paper measures the scaffold on a theory-of-mind benchmark with a
sign-agnostic registration and reports whatever the interval says.

The three outcomes are therefore three readouts of one arrangement. The
asker's pull on a single trace is pillar one, the representation of other
minds that the arrangement requires is pillar two, and the readable dissent
that separate seats produce is pillar three. Two results tie the pillars to
each other rather than merely placing them side by side. Narration does not
create the collective's structure, so internal and external opposition are
different objects with different evidence. And without role lock the
external structure collapses to correlated readers, so external opposition
rests on the control that assignment gives over what a copy will say.
